\documentclass{article}
\usepackage{amsmath, amssymb, amsthm}
\usepackage{hyperref}
\usepackage{geometry}
\geometry{margin=1in}

\title{Erd\H{o}s Problem \#628}
\date{Last updated: 2025-08-31}
\author{Source: erdosproblems.com}

\begin{document}
\maketitle

\noindent\textbf{Status:} FALSIFIABLE \quad \textbf{No prize}

\noindent\textbf{Tags:} graph theory, chromatic number

\medskip
\noindent\textbf{Note:} Erdős-Lovász Tihany Conjecture

\medskip

\noindent\textbf{Problem Statement:}

Let $G$ be a graph with chromatic number $k$ containing no $K_k$. If $a,b\geq 2$ and $a+b=k+1$ then must there exist two disjoint subgraphs of $G$ with chromatic numbers $\geq a$ and $\geq b$ respectively?




This property is sometimes called being $(a,b)$-splittable. A question of Erd\H{o}s and Lov\'{a}sz (often called the Erd\H{o}s-Lov\'{a}sz Tihany conjecture). Erd\H{o}s \cite{Er68b} originally asked about $a=b=3$ which was proved by Brown and Jung \cite{BrJu69} (who in fact prove that $G$ must contain two vertex disjoint odd cycles). Balogh, Kostochka, Prince, and Stiebitz \cite{BKPS09} have proved the full conjecture for quasi-line graphs and graphs with independence number $2$. For more partial results in this direction see the comprehensive survey of this problem by Song \cite{So22}. See also the entry in the graphs problem collection .




References

[BKPS09] Balogh, J\'{o}zsef and Kostochka, Alexandr V. and Prince, Noah and Stiebitz, Michael, The Erd\H{o}s-Lov\'{a}sz Tihany conjecture for quasi-line graphs . Discrete Math. (2009), 3985-3991.

[BrJu69] Brown, W. G. and Jung, H. A., On odd circuits in chromatic graphs . Acta Math. Acad. Sci. Hungar. (1969), 129-134.

[Er68b] Erd\H{o}s, P., Problem 2 . Theory of Graphs (1968), 361.

[So22] Song, Zi-Xia, A survey on the {E}rd\H{o}s-{L}ov\'{a}sz Tihany conjecture . Adv. Math. (China) (2022), 259--274.

\bigskip
\noindent\small{Source: \url{https://www.erdosproblems.com/628}}
\end{document}
